\documentclass[11pt]{article}
\usepackage[utf8]{inputenc}
\usepackage[T1]{fontenc}
\usepackage[italian]{babel}
\usepackage{geometry}
\geometry{a4paper, margin=1in}
\usepackage{xcolor}
\usepackage{listings}
\usepackage{hyperref}

\usepackage{tcolorbox}
\tcbuselibrary{breakable}
\begin{document}

\begin{tcolorbox}[colback=blue!5!white,colframe=blue!75!black,title=\textbf{DOMANDA}, breakable]
\texttt{form-get.html}
\begin{lstlisting}[language=HTML, basicstyle=\ttfamily\small]
<!DOCTYPE html><html>
<body>
<h2>The method Attribute</h2>
<p>This form will be submitted using the GET method:</p>
<form action="/action" target="_blank" method="get">
<label for="fname">First name:</label><br>
<input type="text" id="fname" name="fname" value="John"><br>
<label for="lname">Last name:</label><br>
<input type="text" id="lname" name="lname" value="Doe"><br><br>
<input type="submit" value="Submit">
</form>
<p>After you submit, notice that the form values is visible in the address bar
of the new browser tab.</p>
</body>
</html>
\end{lstlisting}

Esercizio
\begin{itemize}
    \item Eseguire il server web \texttt{serverHTTP.c} e con il browser preferito aprire il file \texttt{form-get.html}
    \item Cosa si vede?
    \item Provare ad analizzare il contenuto della connessione \texttt{TCP} con Wireshark.
    \item Cosa si vede?
\end{itemize}
\end{tcolorbox}

\begin{tcolorbox}[colback=green!5!white,colframe=green!75!black,title=\textbf{RISPOSTA}, breakable]
\textbf{Soluzione Esercizio 23:} Web form GET

\textbf{Operazioni Preliminari}
\begin{enumerate}
    \item È stato copiato il server in linguaggio C e il codice HTML per il modulo web forniti dal professore.
    \item Affinché l'invio del modulo comunichi effettivamente con il server in ascolto in locale, è necessario che il file HTML "\texttt{form-get.html}" punti al server giusto. Se il file viene aperto dal browser come file locale, il campo \texttt{action="/action"} andrebbe modificato in \texttt{action="http://127.0.0.1:8000/action"} (oppure \texttt{http://localhost:8000/action}).
\end{enumerate}

\textbf{Risposta alla Domanda 1: "Cosa si vede (sul browser)?"}\\
Cliccando sul pulsante "Submit", il browser raccoglie i dati inseriti nei campi di input (\texttt{fname} e \texttt{lname}) e li appende direttamente all'URL di destinazione indicato nell'attributo "action", separandoli con il carattere "?" e unendo le varie coppie chiave-valore con la "\&".
Nella barra degli indirizzi della nuova scheda che si apre vedremo quindi una stringa simile a questa:\\
\texttt{http://127.0.0.1:8000/action?fname=John\&lname=Doe}

A schermo si vedrà la pagina statica HTML di cortesia ("Hello World...") restituita dal nostro \texttt{serverHTTP}.

\textbf{Risposta alla Domanda 2: "Cosa si vede (analizzando il TCP con Wireshark)?"}\\
Catturando il traffico sull'interfaccia di loopback, ad esempio tramite il comando:\\
\texttt{sudo tshark -i lo -f "tcp port 8000" -Y "http" -V}\\
oppure seguendo lo stream \texttt{TCP} nella GUI di Wireshark, si osserverà la richiesta \texttt{HTTP} in chiaro, che mostrerà quanto segue:

\begin{lstlisting}[basicstyle=\ttfamily\small]
GET /action?fname=John&lname=Doe HTTP/1.1
Host: 127.0.0.1:8000
User-Agent: ...
\end{lstlisting}

\begin{itemize}
    \item Si nota chiaramente che la particolarità del metodo GET per l'invio dei form consiste nel passare i parametri direttamente nell'URL (Request URI).
    \item Di conseguenza, NON è presente un corpo del messaggio (Body o payload aggiuntivo nella richiesta \texttt{HTTP}).
    \item Questo metodo espone i dati inseriti direttamente nel log delle richieste del server e nei log dei sistemi intermedi o cronologie del browser, rendendolo completamente inadatto al trasferimento di dati sensibili (come le password).
\end{itemize}

\textbf{Comandi Utili per Riprodurre l'Esercizio:}
\begin{lstlisting}[language=bash, basicstyle=\ttfamily\small]
# Eseguire il server
./serverHTTP

# Per testare l'invio via terminale e simulare il submit del form:
curl -v "http://127.0.0.1:8000/action?fname=John&lname=Doe"

# Analizzare l'esatta richiesta arrivata usando TShark:
sudo tshark -i lo -f "tcp port 8000" -z follow,tcp,ascii,0
\end{lstlisting}
\end{tcolorbox}

\end{document}
